% Zone „Das kennst du schon“ – aus bank/koerper/zone.jsonl (zusammenbau v0.1)
\einheitenkopf[zone][Kennst du schon]{Das kennst du schon}
\setcounter{aufgabe}{0}
%% TODO Grundvorstellungs-Aufgabe als erste Hauptnummer der Zone (2.8) fehlt in der Bank

%% TODO Titel: Ich-kann-Satz fehlt in der Bank (Platzhalter: Kettenname) – Nr. 1
%% TODO Anweisung: ein Satz über der ersten Teilaufgabe fehlt in der Bank – Nr. 1
\begin{aufgabe}{Räumliches Vorstellen}
%% TODO Darstellung für schwach fehlt in der Bank (koerper-zone-f1-v1; Abschnitt „Für schwache Schüler“)
\swz{Ein Spielwürfel wird an den Kanten aufgeschnitten und flach ausgebreitet. Aus wie vielen Quadraten besteht das Netz?}{}
%% TODO Darstellung für schwach fehlt in der Bank (koerper-zone-f1-v3; Abschnitt „Für schwache Schüler“)
\swz{Eine Mauer aus Würfeln ist $7$ Würfel lang, $3$ Würfel hoch und $1$ Würfel tief. Wie viele Quadrate siehst du von vorn, wie viele von oben?}{}
\end{aufgabe}

%% TODO Titel: Ich-kann-Satz fehlt in der Bank (Platzhalter: Kettenname) – Nr. 2
%% TODO Anweisung: ein Satz über der ersten Teilaufgabe fehlt in der Bank – Nr. 2
\begin{aufgabe}{Flächeninhalt von Rechteck, Dreieck, Trapez und Kreis (Grundflächen)}
%% TODO Darstellung für schwach fehlt in der Bank (koerper-zone-f2-v1; Abschnitt „Für schwache Schüler“)
\swz{Rechteck $7$ cm lang, $3$ cm breit – Flächeninhalt?}{}
%% TODO Darstellung für schwach fehlt in der Bank (koerper-zone-f2-v3; Abschnitt „Für schwache Schüler“)
\swz{Trapez: parallele Seiten $6$ cm und $3{,}4$ cm, Höhe $5$ cm – Flächeninhalt?}{}
\end{aufgabe}

%% TODO Titel: Ich-kann-Satz fehlt in der Bank (Platzhalter: Kettenname) – Nr. 3
%% TODO Anweisung: ein Satz über der ersten Teilaufgabe fehlt in der Bank – Nr. 3
\begin{aufgabe}{Längen- und Volumeneinheiten, Kubikdezimeter als Liter, Kubikzentimeter als Milliliter, Umrechnen zwischen beiden}
%% TODO Darstellung für schwach fehlt in der Bank (koerper-zone-f3-v1; Abschnitt „Für schwache Schüler“)
\swz{$4$ l – wie viel ml?}{}
%% TODO Darstellung für schwach fehlt in der Bank (koerper-zone-f3-v3; Abschnitt „Für schwache Schüler“)
\swz{$2{,}35$ m – wie viel cm?}{}
\end{aufgabe}

%% TODO Titel: Ich-kann-Satz fehlt in der Bank (Platzhalter: Kettenname) – Nr. 4
%% TODO Anweisung: ein Satz über der ersten Teilaufgabe fehlt in der Bank – Nr. 4
\begin{aufgabe}{Formel nach einer Größe umstellen (V = G · h → h = V : G)}
%% TODO Darstellung für schwach fehlt in der Bank (koerper-zone-f4-v1; Abschnitt „Für schwache Schüler“)
\swz{$P = q \cdot r$ – nach $r$ umstellen?}{}
%% TODO Darstellung für schwach fehlt in der Bank (koerper-zone-f4-v3; Abschnitt „Für schwache Schüler“)
\swz{$A = g \cdot t$ mit $A = 7{,}5$ und $g = 2{,}5$ – $t$?}{}
\end{aufgabe}

%% TODO Titel: Ich-kann-Satz fehlt in der Bank (Platzhalter: Kettenname) – Nr. 5
%% TODO Anweisung: ein Satz über der ersten Teilaufgabe fehlt in der Bank – Nr. 5
\begin{aufgabe}{Multiplizieren mit Dezimalzahlen und π, Runden; Quadrieren und Wurzelziehen}
%% TODO Darstellung für schwach fehlt in der Bank (koerper-zone-f5-v1; Abschnitt „Für schwache Schüler“)
\swz{$0{,}5 \cdot 18$ – wie viel?}{}
%% TODO Darstellung für schwach fehlt in der Bank (koerper-zone-f5-v3; Abschnitt „Für schwache Schüler“)
\swz{$\pi \cdot 6$ – auf eine Stelle gerundet?}{}
\end{aufgabe}

%% TODO Titel: Ich-kann-Satz fehlt in der Bank (Platzhalter: Kettenname) – Nr. 6
%% TODO Anweisung: ein Satz über der ersten Teilaufgabe fehlt in der Bank – Nr. 6
\begin{aufgabe}{Prozentwert berechnen (zehn Prozent von einem Volumen mit Komma)}
%% TODO Darstellung für schwach fehlt in der Bank (koerper-zone-f6-v1; Abschnitt „Für schwache Schüler“)
\swz{$50\,\%$ von $8$ l – wie viel?}{}
%% TODO Darstellung für schwach fehlt in der Bank (koerper-zone-f6-v3; Abschnitt „Für schwache Schüler“)
\swz{$20\,\%$ von $3{,}5$ l – wie viel?}{}
\end{aufgabe}

%% TODO Titel: Ich-kann-Satz fehlt in der Bank (Platzhalter: Kettenname) – Nr. 7
%% TODO Anweisung: ein Satz über der ersten Teilaufgabe fehlt in der Bank – Nr. 7
\begin{aufgabe}{Flächeninhalt von Rechteck, Dreieck, Trapez und Kreis (Grundflächen) – Fehler finden}
\swz[3]{Eine runde Tischplatte hat den Durchmesser $1{,}4$ m. Tom rechnet den Flächeninhalt: \rechnung{\pi \cdot 1{,}4^2 &\approx 6{,}16 \text{ m}^2} Wo ist der Fehler?}{}
\end{aufgabe}

%% TODO Titel: Ich-kann-Satz fehlt in der Bank (Platzhalter: Kettenname) – Nr. 8
%% TODO Anweisung: ein Satz über der ersten Teilaufgabe fehlt in der Bank – Nr. 8
\begin{aufgabe}{Flächeninhalt von Rechteck, Dreieck, Trapez und Kreis (Grundflächen) – selbst rechnen}
%% TODO Darstellung für schwach fehlt in der Bank (koerper-zone-f2-v7; Abschnitt „Für schwache Schüler“)
\swz{Ein runder Teppich hat den Durchmesser $2{,}2$ m – Flächeninhalt auf zwei Stellen?}{}
\end{aufgabe}
